\scanpage{30}
\textbf{C. R. Acad. Sc. Paris, t. 278 (27 mai 1974)}\hfill\textbf{Série A ---}\mspage{1401}\par\noindent\rule{\linewidth}{.35pt}\par

TOPOLOGIE. --- \textit{Classification topologique des plongements des 1-complexes compacts dans les surfaces.} Note (*) de M. \textbf{Yves Ladegaillerie}, présentée par M. Henri Cartan.

{\small On étudie d'abord le cas des surfaces compactes orientées en utilisant une notion d'équivalence de plongements par homéomorphisme. On caractérise alors les classes d'équivalence à l'aide d'invariants.}

1. \textsc{Catégorie des 1-complexes a isotopie près.}

1.1. On appellera \textit{étoile standard} à \(n\) branches (\(n\geqq1\)), la partie de \(\mathbf{R}^2\) définie en coordonnées polaires par \(\{(\rho,\theta),0\leqq\rho\leqq1,\theta=(2\,k\,\pi/n), k\in\mathbf{Z}\}\). On appellera \textit{étoile} tout espace topologique homéomorphe à une étoile standard ou réduit à un point (0 branche). Un \textit{1-complexe compact} est un espace topologique compact dont tout point a un voisinage homéomorphe à une étoile à \(n\) branches; si \(n\ne2\), on dira que c'est un \textit{sommet}, si \(n=0\) c'est un \textit{sommet isolé}.

1.2. On notera \(\mathscr{K}_{\mathscr I}\) la catégorie des \textit{1-complexes à isotopie près}, dont les objets son\textsuperscript{t} les 1-complexes compacts et les morphismes sont les \textit{classes d'isotopie} d'homéomorphismes. On notera \(\mathscr{K}_{\mathscr C}\) la catégorie des \textit{1-complexes combinatoires} dont les objets sont les 5-uples \((\mathrm{S},\mathrm{A},\mathrm{A}_1,\sigma,s)\), où \(\mathrm{S},\mathrm{A},\mathrm{A}_1\) sont des ensembles finis, \(\mathrm{A}_1\subset\mathrm{A}\), \(\sigma\) une involution sans point fixe de \(\mathrm{A}\) telle que \(\sigma(\mathrm{A}_1)=\mathrm{A}_1\), \(s\) une application de \(\mathrm{A}_1\) dans \(\mathrm{S}\) telle qu'il n'existe pas de \(x\) de \(\mathrm{S}\) tel que \(\operatorname{card}(s^{-1}(x))=2\). Un morphisme entre \((\mathrm{S},\mathrm{A},\mathrm{A}_1,\sigma,s)\) et \((\mathrm{S}',\mathrm{A}',\mathrm{A}'_1,\sigma',s')\) est un couple \((\varphi,\psi)\), où \(\varphi\) est une bijection de \(\mathrm{S}\) sur \(\mathrm{S}'\), \(\psi\) une bijection de \(\mathrm{A}\) sur \(\mathrm{A}'\) telles que \(\psi(\mathrm{A}_1)=\mathrm{A}'_1\), \(\varphi\circ s=s'\circ\psi\vert_{\mathrm{A}_1}\), \(\psi\circ\sigma=\sigma'\circ\psi\).

1.3. On définit un foncteur \(\mathscr{C}\) de \(\mathscr{K}_{\mathscr I}\) dans \(\mathscr{K}_{\mathscr C}\) en associant au 1-complexe \(\mathrm{K}\) le 5-uple \((\mathrm{S},\mathrm{A},\mathrm{A}_1,\sigma,s)\), où \(\mathrm{S}\) est l'ensemble (fini) des sommets de \(\mathrm{K}\), \(\mathrm{A}\) l'ensemble (fini) des \textit{arcs} de \(\mathrm{K}\) (un arc étant une arête orientée et une arête une composante connexe de \(\mathrm{K}\setminus\mathrm{S}\)), \(\mathrm{A}_1\) l'ensemble des arcs ayant une source, \(s\) l'application de \(\mathrm{A}_1\) dans \(\mathrm{S}\) qui à un tel arc associe sa source, \(\sigma\) l'involution de \(\mathrm{A}\) qui a un arc associe son inverse.

1.4. \textsc{Théorème.} --- \textit{Le foncteur \(\mathscr{C}\) est une équivalence de catégories.}

2. \textsc{Préliminaires topologiques.} --- \textit{Terminologie.} --- Une surface (resp. surface à bord) est une 2-variété topologique (resp. 2-variété à bord).

2.1. On appellera \textit{plongement} d'un 1-complexe \(\mathrm{K}\) dans une surface à bord orientée \(\mathrm{V}\) : une injection continue de \(\mathrm{K}\) dans l'intérieur de \(\mathrm{V}\).

2.2. \textsc{Proposition.} --- \textit{Toute injection continue de \([0,1]\) dans \(\mathbf{R}^2\) se prolonge en un homéomorphisme de \(\mathbf{R}^2\) sur lui-même} [\textit{voir} (\textsuperscript{2}), (\textsuperscript{3}), (\textsuperscript{4})].

2.3. \textsc{Proposition.} --- \textit{Tout plongement d'une étoile \(\mathrm{K}\) dans \(\mathbf{R}^2\) orienté détermine de façon canonique un ordre circulaire sur l'ensemble des branches de \(\mathrm{K}\) et est isotope par isotopie ambiante à un plongement standard} (1.1) \textit{induisant le même ordre circulaire.}

2.4. \textsc{Définition.} --- Deux plongements \(f_1\) et \(f_2\) d'un 1-complexe compact \(\mathrm{K}\) dans deux surfaces orientées à bords \(\mathrm{V}_1\) et \(\mathrm{V}_2\) sont dits \textit{équivalents} s'il existe un homéomorphisme orienté \(h\) de \(\mathrm{V}_1\) sur \(\mathrm{V}_2\) tel que \(h\circ f_1=f_2\).
% archive sheet number: 29

